\section{The finite certificate at \texorpdfstring{$m=2$}{m=2}}
\label{sec:certificate}

Fix even indices $2\leq i<j\leq364$, and write
\[
 v_k=(1,2^k,3^k)^{\mathsf T},\qquad
 q(2)=(q_1(2),q_2(2),q_3(2))^{\mathsf T}.
\]
The integer vector $w=v_i\times v_j$ is perpendicular to both columns,
so the combination $\sum_{b=1}^3w_br_{n,b}$ cancels $\beta(i)$ and
$\beta(j)$ for every $n$. Its constant term at $n=2$ is
\[
 q_w(2)=w\cdot q(2)=\det\bigl(q(2),v_i,v_j\bigr).
\]
The determinant is the constant left after canceling the two beta
terms. If it is nonzero, Proposition~\ref{prop:dilation} proves
nonvanishing at $n=2p$ for all sufficiently large $p$ under the
rationality assumption. We therefore need the following finite statement.

\begin{lemma}[Exact seed certificate]\label{lem:seed}
For every pair of distinct even indices $2\leq i<j\leq364$,
\begin{equation}\label{eq:determinant}
 \det\begin{pmatrix}
 q_1(2)&1&1\\q_2(2)&2^i&2^j\\q_3(2)&3^i&3^j
 \end{pmatrix}\ne0.
\end{equation}
\end{lemma}

We prove this by reduction modulo the prime $\ell=1\,000\,003$.
A nonzero residue proves that the rational determinant is nonzero,
provided its entries have denominators prime to $\ell$. We first give
the three seed residues and test all pairs. Their derivation follows in
Subsection~\ref{subsec:seed-residues}.

\subsection{The determinant test}

The common normalization $C_2=2^{1825}3^{473}$ is nonzero modulo $\ell$,
so we may compute with
\[
 U_2(t)=\frac{P_2(t)}{t^{365}(t+4)^{365}(t+8)^{365}}.
\]
Let $\widetilde a_{i,h}$ be the partial-fraction coefficients of $U_2$
and $\widetilde q_b$ the corresponding constant terms, defined as in
\eqref{eq:partials} and \eqref{eq:qb}. Then $q_b(2)=C_2\widetilde q_b$.
The denominators of $\widetilde a_{i,h}$ have only the prime factors
$2$ and $3$; the finite sums add only primes at most $23$.
Reduction modulo $\ell$ is therefore legitimate.

The calculation in Subsection~\ref{subsec:seed-residues} gives
\begin{equation}\label{eq:seed-residues}
 (\widetilde q_1,\widetilde q_2,\widetilde q_3)
        \equiv(614660,346158,44336)\pmod\ell.
\end{equation}
For $Q=(614660,346158,44336)$, compute, modulo $\ell$,
\begin{align*}
 D_{i,j}={}&Q_1(2^i3^j-2^j3^i)
          -Q_2(3^j-3^i)+Q_3(2^j-2^i),\\
 &\hspace{32mm}2\leq i<j\leq364,\quad i,j\text{ even}.
\end{align*}
\begin{samepage}
The following loop checks all pairs using the residues
in \eqref{eq:seed-residues}:
\begin{quote}\small
\begin{verbatim}
ell = 1000003
Q1, Q2, Q3 = 614660, 346158, 44336
count = 0
for i in range(2, 365, 2):
    for j in range(i + 2, 365, 2):
        det = (Q1 * (pow(2,i,ell)*pow(3,j,ell)
                     - pow(2,j,ell)*pow(3,i,ell))
               - Q2 * (pow(3,j,ell)-pow(3,i,ell))
               + Q3 * (pow(2,j,ell)-pow(2,i,ell))) % ell
        assert det != 0
        count += 1
assert count == 16471
\end{verbatim}
\end{quote}
\end{samepage}
All $\binom{182}{2}=16471$ residues are nonzero. Since $C_2$ is also
nonzero modulo $\ell$, this proves \eqref{eq:determinant} over $\Q$.
The weights $(15,-6,1)$ from the introduction give
\begin{equation}\label{eq:one-pair}
 15\cdot614660-6\cdot346158+44336\equiv187267\not\equiv0\pmod\ell.
\end{equation}

\subsection{Recovering the seed residues}\label{subsec:seed-residues}

To obtain \eqref{eq:seed-residues}, we compute the partial fractions of
$U_2$. At a pole $-4h$, substitute $t=u-4h$ and multiply by $u^{364}$.
The resulting function is regular at $u=0$. Its first $364$ Taylor
coefficients give the principal part of $U_2$, in reverse order.

At the pole $-4h$, where $h=0,1,2$, form the multiset
\[
 \mathcal C_h=\{-4h-14\}\mathbin{\cup}
 \{\!-4h-14+j+\theta:0\leq j<36,\ \theta\in\Theta\}.
\]
It has exactly one zero. Delete that zero to obtain $\mathcal C_h^*$,
and put $\mathcal D_h=\{4(k-h):0\leq k\leq2,\ k\ne h\}$.
Thus
\[
 B_h(u):=u^{364}U_2(u-4h)
       =\frac{\prod_{c\in\mathcal C_h^*}(u+c)}
                    {\prod_{d\in\mathcal D_h}(u+d)^{365}}
       =\sum_{j\geq0}\gamma_{h,j}u^j.
\]
Its logarithmic derivative gives a convenient exact recurrence:
\begin{align}
 \gamma_{h,0}&=\frac{\prod_{c\in\mathcal C_h^*}c}
                         {\prod_{d\in\mathcal D_h}d^{365}},\notag\\
 T_{h,k}&=\sum_{c\in\mathcal C_h^*}c^{-k}
                      -365\sum_{d\in\mathcal D_h}d^{-k},\notag\\
 \gamma_{h,j}&=\frac1j\sum_{k=1}^{j}(-1)^{k-1}T_{h,k}\gamma_{h,j-k}
                 \quad(1\leq j\leq363),\label{eq:seed-recurrence}\\
 \widetilde a_{i,h}&=\gamma_{h,364-i}\quad(1\leq i\leq364).
 \notag
\end{align}
All divisions are valid modulo $\ell$. Each nonzero
$c\in\mathcal C_h^*$ and $d\in\mathcal D_h$ has numerator and denominator
of absolute value less than $\ell$, and each recurrence index $j$ is
at most $363$.

Substitution in \eqref{eq:qb} gives the residues in
\eqref{eq:seed-residues}. The supplied implementations in
\texttt{checks/seed.py} and \texttt{checks/verify.py} use only Python's
standard library. Primality of $\ell$ is checked by trial division up to
its integer square root.

The checker computes all $1092$ partial-fraction coefficients both by
\eqref{eq:seed-recurrence} and by multiplying truncated power series.
The results agree in every entry. It also checks rational-function
reconstruction at three points and the vanishing of all odd $\rho_i$
modulo $\ell$. A separate calculation over $\Q$ verifies the
divisibility assertions of Proposition~\ref{prop:denominators} at
$n=2$ for all $1092$ coefficients and all three constant terms, and
recovers \eqref{eq:seed-residues}.

